\subsection{The MANO Forward Model}\label{subsec:forward_model}

The \texttt{KinematicModel} class in \texttt{models} realises the \acrshort{mano}
forward pass: given pose and shape parameters it produces a deformed mesh and the
joint and fingertip keypoints, following the linear-blend-skinning formulation of
\autoref{subsec:mano}. On construction it loads the arrays produced by the offline
conversion --- the pose \acrshort{pca} basis and mean, the joint regressor, the
skinning weights, the pose and shape blend-shape bases, the template mesh, the
face topology, and the parent indices.

Pose may be supplied in either of two forms through \texttt{set\_params}: as a set
of \acrshort{pca} coefficients, or as absolute axis-angle rotations. The
inverse-kinematics path of this system uses the latter, with the global rotation
occupying the first slot of the pose vector. The \texttt{update} routine then
recomputes the mesh. It first applies the shape blend shapes to the template and
regresses the joint positions from the resulting vertices. Each joint's axis-angle
rotation is converted to a rotation matrix by Rodrigues' formula, and the global
transforms are composed down the kinematic tree, each joint's transform built on
that of its parent. The rest-pose transform is then subtracted so that the zero
pose leaves the mesh undeformed, and linear blend skinning combines the
per-joint transforms, weighted by the skinning weights, to produce the final
posed vertices.

% One detail of the implementation is worth stating precisely against the
% formulation in \autoref{subsec:mano}. Although the pose blend-shape basis
% $B_P$~(\autoref{eq:mano_bp}) is loaded, it is not applied in \texttt{update}: the
% runtime mesh is the shape-adjusted template deformed by skinning alone, and the
% pose-corrective term that \acrshort{mano} adds to reduce joint-collapse artefacts
% is omitted. The shape contribution $B_S$ of~\autoref{eq:mano_tp} is retained.

Finally, the keypoints are read from the posed mesh through an extended joint
regressor, which augments the standard regressor with five additional rows that
each select a single fingertip mesh vertex. This yields the 21-keypoint layout
that matches the detector's output and that the inverse-kinematics solver expects.
